\section{Introduction}\label{sec:introduction}

One-relator presentations impose a single equation on a free group, but their geometric behavior ranges from free groups to groups with distorted cyclic subgroups. Hyperbolicity rules out the latter behavior. The purpose of this paper is to show that it also forces a strong finite-cover geometry.

\begin{theorem}\label{thm:main}
Let $X$ be finite, let $r\in F(X)$, and suppose that
\[
G=F(X)/\langle\!\langle r\rangle\!\rangle
\]
is word-hyperbolic. Then $G$ is virtually compact special. More precisely, there are a finite-index subgroup $H\leq G$, a finite connected nonpositively curved cube complex $C$, and a finite simplicial graph $\Lambda$ such that
\[
\pi_1(C)\cong H
\quad\text{and}\quad
C\longrightarrow S_\Lambda
\]
is a cubical local isometry to the Salvetti complex of the right-angled Artin group defined by $\Lambda$.
\end{theorem}

The compactness assertion is part of the conclusion. An abstract embedding of a finite-index subgroup in a right-angled Artin group would not by itself supply the model required in Theorem~\ref{thm:main}.

The inductive study of one-relator groups begins with Magnus's
Freiheitssatz: omitting a generator that occurs in a cyclically reduced
relator leaves a subgroup freely generated by the remaining generators
\cite{Magnus1930}. Magnus subsequently solved the word problem for
one-relator groups \cite{Magnus1932}. These results established the
importance of reducing a one-relator problem to simpler presentations
while retaining control of distinguished free subgroups. The Magnus--Moldavanskii
hierarchy develops this approach through splittings over those subgroups
\cite[Sections~1.1--1.3]{LintonHierarchies2025}.
The geometric question is whether the splittings have the quasiconvexity
needed to construct finite special covers.

A second line of inquiry concerns finite quotients. A group is
\emph{residually finite} if each nonidentity element survives in some
finite quotient. Baumslag's 1967 Conjecture~A asks for this property for
one-relator groups with nontrivial torsion \cite{Baumslag1967Residual}.
His paper contrasts this question with non-Hopfian one-relator examples,
which cannot be residually finite. Wise's quasiconvex hierarchy theory
proved the stronger virtual compact specialness conclusion for
one-relator groups with torsion \cite[Corollary~19.2]{Wise2021}.
Theorem~\ref{thm:main} extends the compact-special conclusion, and hence
residual finiteness, to every hyperbolic one-relator group.

The cubical approach explains why this strengthening matters.
Haglund--Wise's special cube complexes connect nonpositive curvature
with embeddings in right-angled Artin groups and, in the hyperbolic
case, separation of quasiconvex subgroups in finite quotients \cite{HaglundWise2008}.
Wise's hierarchy theorem makes quasiconvex splittings a way to obtain
this structure \cite[Theorem~13.1]{Wise2021}; Agol's theorem supplies
it for hyperbolic groups acting properly and cocompactly on CAT(0)
cube complexes \cite[Theorem~1.1]{Agol2013}.
Thus the geometric conclusion controls both linear representations and
finite-quotient information. Section~\ref{sec:free-by-cyclic} explains
these consequences, including conjugacy separation and virtual
retractions, as well as the fibering conclusions below.

In the torsion-free direction, Linton proved hyperbolicity and virtual
specialness when all two-generator subgroups of a one-relator group are
free \cite[Theorem~7.2 and Section~1]{LintonHierarchies2025}.
His subsequent work identifies primitive extension groups as the
terminal cases in the general hyperbolic problem and specifies the
quasiconvex towers through which a terminal result propagates
\cite{LintonHyperbolic2026}. The remaining task is therefore a
geometric construction for those primitive groups. Here we construct
virtually compact special quotients, recover that structure using
selective covers of the fillings, and then propagate it up the truncated
tower. Wise singled out hyperbolic one-relator groups in his broader
question about cocompact cubulation of hyperbolic groups with
hierarchies \cite[Problem~13.15]{Wise2014Cubical}.
Theorem~\ref{thm:main} establishes the compact-special conclusion for
this one-relator class.

We use \emph{free-by-cyclic} in the subgroup-closed sense: a group $H$
fits into an exact sequence
\[
1\longrightarrow K\longrightarrow H\longrightarrow Q\longrightarrow1
\]
with $K$ free of possibly infinite rank and $Q\leq\mathbb Z$, possibly
trivial. Wise's Conjecture~17.8 asks whether every hyperbolic one-relator
group is virtually free-by-cyclic, as recalled by Kielak--Linton
\cite[Section~1.1]{KielakLinton2024Virtually}. Their Corollary~1.3
identifies virtual compact specialness as a sufficient additional
hypothesis. Combining that corollary with Theorem~\ref{thm:main} gives a positive
answer. Corollary~\ref{cor:free-by-cyclic} also records the
Baumslag--Solitar-free consequence from the companion hyperbolicity
theorem and a stronger finite-index embedding into a free-by-$\mathbb Z$
group whose free kernel is finitely generated.

In the proof we use the hyperbolicity and Magnus-subgroup results of
\cite{OpenAIHyperbolicity} at two specific points: embedded primitive
extension groups are themselves hyperbolic, and their actual Magnus
rows have the required hyperbolic geometry. The criterion in that paper
excludes every Baumslag--Solitar subgroup $\mathrm{BS}(m,n)$ with
$mn\ne0$. No virtual-specialness conclusion is imported from it.

The recovery step is a filling criterion proved in
Section~\ref{sec:01-filling}. Its input is a coherent, locally indicable
hyperbolic group of cohomological dimension at most two, with a finite
classifying complex acyclic over its Hughes-free division field, and
arbitrarily deep virtually compact special fillings along specified
peripherals. We replace these by \emph{selective} fillings that retain
any prescribed finite collection of kernel factors. Their covers have
a quasiconvex cubical hierarchy. This geometric control realizes the
finite coefficient data of a chain contraction in ordinary Novikov
completions for both signs of one character. The resulting
finite-index subgroup maps onto $\mathbb Z$ with a \emph{finite-rank}
free kernel, and hence is virtually compact special.

Before constructing the primitive fillings, we use this criterion to
prove a second result. It will also supply the specialness of the
periodic end groups encountered in that construction.

\begin{theorem}\label{thm:ascending-intro}
Let $F$ be a finitely generated free group and let $\phi:F\to F$ be injective. If its ascending mapping torus
\[
F*_{\phi}=\langle F,t\mid t^{-1}ft=\phi(f),\ f\in F\rangle
\]
is word-hyperbolic, then it is virtually compact special.
\end{theorem}

Hagen and Wise cubulate hyperbolic free-by-cyclic groups with finite-rank
free kernel \cite[Theorem~A]{HagenWise2016}. Their irreducible theorem also covers expanding,
fundamental-group-injective train-track maps whose transition matrices
are irreducible and whose mapping-torus groups are hyperbolic
\cite[Theorem~B]{HagenWise2015}. They posed the cubulation
problem for hyperbolic ascending HNN extensions defined by general
monomorphisms in \cite[Problem~C]{HagenWise2016}.
Mutanguha constructs expanding relative immersions for injective
nonsurjective endomorphisms and proves the
Baumslag--Solitar criterion for arbitrary injective ascending tori
\cite[Theorems~4.5 and~6.7]{Mutanguha2021Dynamics}.
He also records Kielak's question whether hyperbolic ascending tori of
injective nonsurjective free-group endomorphisms are virtually
free-by-cyclic \cite[Overview, p.~3, Question (Virtual fibering)]{Mutanguha2021Dynamics}.
The proof below retains arbitrary injective input and establishes the
compact-special conclusion. Linton proved that every
hyperbolic free-by-cyclic group is virtually compact special, allowing
its free kernel to have infinite rank
\cite[Corollary~1.5]{Linton2026Embedding}. The canonical kernel of a
strict ascending torus is presented here as an increasing union of
free groups; freeness of that union is not an input to our argument.

After Theorem~\ref{thm:ascending-intro} is proved, the ascending group's
field acyclicity and virtual RFRS property give a stronger fibering
conclusion: it has a finite-index subgroup mapping onto $\mathbb Z$
with a finite-rank free kernel
(Corollary~\ref{cor:ascending-finite-rank}). The resulting character is
not identified with the canonical ascending height. For general one-relator groups we
retain the possibly infinite-rank kernel in
Corollary~\ref{cor:free-by-cyclic}.

Kielak's virtual-fibering method realizes division-field coefficients
in ordinary Novikov completions along an RFRS filtration
\cite[Section~4.2 and Theorem~5.2]{Kielak2020}. Fisher's later work
relates this method to cohomological dimension of kernels
\cite{Fisher2026Kernels}. Here the original group is not yet known to
be virtually RFRS: the selective switch supplies the required finite
coefficient data before the two-sided Novikov argument. Once compact
specialness is established, Kielak--Linton's theorems give the general
one-relator fibering and ambient-embedding consequences
\cite{KielakLinton2024Virtually}.

The remaining construction provides the primitive group's good
fillings. Its Magnus splitting has a height coordinate. We call the
vertex and edge spaces at specified heights its \emph{rows}; these are
the actual spaces in the developed Bass--Serre tree, not merely their
quotient labels. Long paths with noncyclic stabilizers obstruct an
acylindrical hierarchy. Section~\ref{sec:03-rows} controls them by a
joint injection into one finite-dimensional homology space. Keeping
this injection under subgroup restriction, rather than only its rank
inequality, distinguishes coincident contributions from different
path orbits.

Expanding trajectories then identify periodic ends and join them
along bilateral trajectories (Section~\ref{sec:05-gates}). Homology
prevents circuits and rank growth in the resulting incidence trees.
We close the corresponding quasiconvex groups under cyclic
intersections, fold rows, and enlarge shared ends until they form a
malnormal peripheral system. There are two termination arguments:
horizontal joins with no homological growth stabilize because their
increasing union is finitely generated, whereas each genuine
enlargement at a shared line end strictly increases the injected
homology dimension. The
coefficient constructions in Sections~\ref{sec:04-coefficients}
and~\ref{sec:analytic-localization} preserve the joint injections
through these operations. The first retains prescribed subgroup
supports in scalar chain equations; the second localizes semilinear
affiliated coefficients to finitely many periodic cosets.

We fill finite-index subgroups of the peripheral \emph{height
kernels}, retaining the height direction. Exact intersection transfer
(Section~\ref{sec:transfer-relative}) makes the resulting splitting
acylindrical. The compatible kernels of Section~\ref{sec:08-kernels}
also give exact link covers for the filled rows. Binary labels on
these links have sparse changes; the changes specify horizontal wall
attachments. Within each filled vertex row, panels join through
midpoint slots in the central pieces and then through satellite and
outer fans. This produces quasiconvex trees carrying binary cocycles
(Sections~\ref{sec:09-fans} and~\ref{sec:10-boundary}). Keeping central
and satellite panels separate lets the labels control departures
introduced by later attachments, including at root extensions.

Three boundary configurations require different panel tests. The
resulting separators satisfy Bergeron--Wise's boundary criterion
\cite[Theorem~1.4]{BergeronWise2012}; Agol's theorem then gives virtual
compact specialness of the filled rows. The transferred hierarchy
assembles the good quotients in Theorem~\ref{boundary:good-fillings}.
Finally Theorem~\ref{boundary:primitive-vcs} applies the filling
criterion to the unfilled primitive group, and the truncated
one-relator tower proves Theorem~\ref{thm:main}.

\section{Background, conventions, and the reduction}\label{sec:background}

We write VCS for \emph{virtually compact special}, QC for quasiconvex, and fg for finitely generated. RFRS means residually finite rationally solvable. An ascending mapping torus always uses an injective endomorphism unless explicitly described as a presentation torus before passing to its stable image. All subgroup occurrences in a Bass--Serre tree include their actual conjugating frames. Thus two occurrences of isomorphic subgroups need not be the same subgroup or the same peripheral object.

\subsection{Cubulation and hierarchy results}

We use the following standard results in their compact forms.

\begin{proposition}[Cubulation and quasiconvex hierarchies]\label{prop:compact-toolkit}
\begin{enumerate}
\item A word-hyperbolic group acting properly and cocompactly on a CAT(0) cube complex is virtually compact special.
\item A hyperbolic group with a finite graph-of-groups decomposition whose edge groups are quasiconvex in the ambient group and whose vertex groups are hyperbolic VCS is VCS. The same applies to a finite hierarchy of such decompositions.
\item A compact special cube complex admits a cubical local isometry to a Salvetti complex. Its defining graph can be taken finite.
\end{enumerate}
\end{proposition}

These are the virtual special theorem \cite[Theorem~1.1]{Agol2013}, the quasiconvex hierarchy theorem \cite[Theorem~13.1]{Wise2021}, and the special-complex characterization \cite[Proposition~3.10 and Theorem~4.2]{HaglundWise2008}, respectively. For the last assertion, only finitely many hyperplanes and crossing relations occur in a finite complex. Passing to a further special cover, when needed, provides the oriented two-sided convention.

The hierarchy statement with VCS terminal groups follows also by cutting a compact special cover of each terminal model along its finitely many two-sided hyperplane families. The cutting traces are convex, hence quasiconvex, and the resulting hierarchy ends in trivial groups. Finite-index supergroups are allowed in the hyperbolic hierarchy class. In any finite splitting used here, quasiconvexity of an edge group is asserted in the ambient group, not merely in an adjacent vertex group. Finitely generated vertex groups are then undistorted: in a tree of Cayley spaces, replace excursions from a target vertex through incident edge cosets by paths in those cosets of comparable length.

We use residual finiteness, virtual RFRS, quasiconvex subgroup separability and virtual retractions for hyperbolic VCS groups \cite{Agol2008,HaglundWise2008,HaglundWise2012,Wise2021}. The malnormal special quotient theorem supplies designated finite-index normal subgroups $\dot P_i\unlhd P_i$ for a finite almost malnormal quasiconvex peripheral family. Normal finite-index kernels $N_i\unlhd P_i$ contained in $\dot P_i$ give hyperbolic VCS quotients \cite[Theorem~2.7 and Corollary~2.8]{AGM2016}. We retain this containment condition at every use; it is not replaced by finite avoidance. Additional finite avoidance conditions can be imposed by further intersections. Exact peripheral embeddings follow from the relative Dehn filling theorem \cite[Theorem~1.1]{Osin2007}.

For relative hyperbolicity we use peripheralization of finite almost malnormal quasiconvex families and the uniform relative Dehn filling theorem \cite{Bowditch2012,Osin2006,Osin2007}. For rows that are only relatively finitely generated, the fine-graph to relative-presentation comparison is \cite[Proposition~3.47]{GerasimovPotyagailo2015}; its hypotheses and exact infinite vertex stabilizers are checked at the constructions below. In particular, bounded isolated peripheral components in sufficiently long fillings have labels from a fixed finite list in the \emph{original} peripheral groups. Later arguments explain the additional tests needed for exact subgroup-intersection transfer; that stronger conclusion is not assumed as part of the ordinary filling theorem.

\subsection{The primitive extension reduction}

Here is the shape of the groups that remain after the reduction. For
$k\ge1$, write $a_i=t^{-i}at^i$. Their presentations have the forms
\[
 \langle a,t\mid\operatorname{pr}_{p/q}(x,y)\rangle,
 \qquad
 \langle a,t\mid\operatorname{pr}_{p/q}(xy,z)\rangle,
\]
where $p,q$ are relatively prime positive integers. The notation
$\operatorname{pr}_{p/q}(v,w)$ means substitution of $v,w$ into the
positive primitive word of slope $p/q$ in two formal letters. Primitivity
means membership in a free basis before substitution. The words $x,y$
use the overlapping windows $a_0,\ldots,a_{k-1}$ and
$a_1,\ldots,a_k$, respectively, and each lies outside their common
subgroup $F(a_1,\ldots,a_{k-1})$. The second form also uses a nontrivial
word $z$ in that common subgroup. These shapes are only a preview:
the slot and primitive exceptional-property conditions in
Definition~\ref{row:primitive-data}, including its malnormality and
exceptional-factorization qualifications, are part of the definition.

\begin{proposition}[Reduction]\label{prop:primitive-reduction}
To prove Theorem~\ref{thm:main}, it suffices to prove that every torsion-free hyperbolic primitive extension group in the two forms of \cite[Section~5]{LintonHyperbolic2026} is VCS.
\end{proposition}

\begin{proof}
The trivial relator gives a free group. Wise's virtual compact specialness theorem handles the torsion case \cite[Corollary~19.2]{Wise2021}. For a torsion-free group, use the maximal factored one-relator tower and primitive extension reduction of \cite[Theorem~4.8, Corollary~4.9, Lemma~5.2 and Theorem~5.6]{LintonHyperbolic2026}. The primitive terminal groups have finite one-relator presentations and embed in $G$.

A hyperbolic group contains no $\mathrm{BS}(m,n)$ with $mn\ne0$: cyclic subgroups generated by infinite-order elements are undistorted and have virtually cyclic commensurators, whereas the indicated Baumslag--Solitar groups violate these alternatives. The embedded primitive groups are therefore BS-free. Theorem~1.1 of \cite{OpenAIHyperbolicity} makes each of them hyperbolic; no general assertion that finitely generated subgroups of hyperbolic groups are hyperbolic is being made.

With these terminal groups hyperbolic, the tower truncated at its primitive extension complex has the required acylindricity and quasiconvex inclusions by \cite[Corollary~4.9]{LintonHyperbolic2026}. Only this truncated tower is used; no acylindricity is asserted for the unused remainder of the maximal tower below it. If each primitive terminal group is VCS, Proposition~\ref{prop:compact-toolkit} propagates VCS through these upward steps. If the tower reaches an elementary terminal group without a primitive obstruction, the same argument starts there. This proves the reduction with the compact conclusion.
\end{proof}

The precise primitive presentations and their actual refined rows are given in Section~\ref{sec:03-rows}. They have two generators and one nontrivial relator, so their standard aspherical complexes have Euler characteristic zero. Torsion-free one-relator groups are locally indicable \cite{Brodskii1984,Howie1982}; see also the direct proof in \cite[Theorem~3.1 and Corollary~3.2]{Howie2000}. They are coherent \cite[Theorem~1.1(1)]{JaikinLinton2025}. For the nontrivial non-power presentations here, Lyndon's asphericity theorem \cite{Lyndon1950} is also supplied by \cite[Corollary~6.5]{AntolinDicksLinnell2010}: specialize to the free group on its Cayley tree and a singleton non-power relator. The root quotient is trivial, so the resulting exact augmented sequence is the cellular sequence of the simply connected presentation cover; that cover is acyclic and hence contractible. The rational Linnell division fields are available, with Hughes-free subgroup compatibility, by \cite{JaikinLopezAlvarez2020,JaikinUniversality2021}.

For such a primitive group, the universal-cover chain complex is acyclic over its division field. The Fox 2-boundary is a nonzero column, hence injective over the field; degree-zero homology vanishes because the group is nontrivial, and Euler characteristic zero then gives degree-one acyclicity. These are the finiteness and field hypotheses used by the filling criterion.

\subsection{Coefficient and tree conventions}

For a subgroup $U$ of the torsion-free ambient group, $D_U$ denotes its rational Hughes-free Linnell field in the affiliated operator algebra $\mathcal U(U)$. Inclusions use the compatible subgroup-field embeddings. Distinct subgroup cosets are independent over the corresponding field. We compute chains with coefficients on the left and differentials given by right matrices; the Fox rule is
\[
d(rs)=d(r)+r\,d(s),\qquad \partial d(r)=r-1.
\]
Changing a subgroup frame by $s$ conjugates group symbols and right-multiplies their coefficients by $s^{-1}$. This convention preserves their ambient homology images. The notation $H_*(V;D_U)$ for $V\le U$ means scalar extension from the intrinsic subgroup field. We use the Bass--Serre homology and cohomology sequences, including infinite quotient graphs: direct sums occur in homology and products in cohomology \cite{Serre1980,Brown1982}.

All matrix traces in the affiliated-operator argument are compatible \emph{unnormalized} extensions of the group trace. This ensures equality of traces for rectangular products in the two possible orders. Matrix dimensions remain finite, so this convention does not change tightness in the measure topology.

\begin{proof}[Proof of Theorem~\ref{thm:main}, assuming the later constructions]
Theorem~\ref{boundary:primitive-vcs} proves the primitive case required by Proposition~\ref{prop:primitive-reduction}. Apply that proposition and then the Salvetti characterization in Proposition~\ref{prop:compact-toolkit}. This gives a finite connected special cover, its finite-index fundamental group, and the stated local isometry. Theorem~\ref{thm:ascending-intro} is established in Section~\ref{sec:02-ascending}.
\end{proof}
